% Generated draft: requires mathematical and source review.
\documentclass[11pt]{article}
\usepackage[margin=25mm]{geometry}
\usepackage{fontspec}
\setmainfont{Latin Modern Roman}
\usepackage{amsmath,amssymb,mathtools,unicode-math}
\setmathfont{Latin Modern Math}
\usepackage{enumitem,needspace,xurl,hyperref}
\setlength{\parindent}{0pt}
\setlength{\parskip}{5pt}
\setlength{\emergencystretch}{3em}
\newcommand{\N}{\mathbb{N}}
\newcommand{\Z}{\mathbb{Z}}
\newcommand{\Q}{\mathbb{Q}}
\newcommand{\R}{\mathbb{R}}
\newcommand{\C}{\mathbb{C}}
\newcommand{\F}{\mathbb{F}}
\newcommand{\A}{\mathbb{A}}
\newcommand{\PP}{\mathbb{P}}
\newcommand{\E}{\mathbb{E}}
\newcommand{\eps}{\varepsilon}
\newcommand{\dd}{\,\mathrm{d}}
\begin{document}
\section*{UnsolvedMath ID: 30000823}
\textbf{Source problem code:} catalogue entry
\subsection{Definitions and mathematical statement}
Let \(M\) be a compact oriented integral homology three-sphere, so \(H_*(M;\Z)\cong H_*(S^3;\Z)\), and let \(L\subset M\) be a link. Its exterior and link group are
\[
 X_L=M\setminus\operatorname{int}\nu(L),\qquad
 G_L=\pi_1(X_L),
\]
where \(\nu(L)\) is a union of disjoint tubular neighborhoods of the components. The exterior and the open complement have the same fundamental group. A graph link, as defined by the source, is one whose exterior has only Seifert-fibered geometric pieces; the decomposition and any irreducibility conventions are those intended in the report.

A group \(\Gamma\) is quasi-projective if there exists a connected smooth complex quasi-projective algebraic variety \(Y\) with \(\Gamma\cong\pi_1(Y)\). The source calls this property global and states the equivalence with quasi-projectivity. It does not require \(Y\) to be projective, require \(X_L\) itself to be algebraic, or require a specified dimension of \(Y\).

The problem asks to determine precisely which non-Seifert graph links satisfy
\[
 \exists\,Y\text{ connected smooth complex quasi-projective},\qquad
 \pi_1(X_L)\cong\pi_1(Y).
\]
Seifert links are excluded from the requested classification. Under the customary interpretation these are links compatible with a Seifert fibration, and their exteriors are Seifert-fibered; the exact report convention, including exceptional exteriors, needs confirmation. Having every geometric piece Seifert-fibered is not the same condition as requiring the whole exterior to have one Seifert fibration.

The supplied assessment dated 2026-08-21 reports broad-family results from splice diagrams and Alexander polynomials, including non-globality for many non-Seifert graph links, but says the unrestricted classification is incomplete. It gives no universal obstruction, complete splice-diagram criterion, or list that could replace the full question.
\subsection{Short English statement}
Among graph links outside the Seifert-link case, which exterior groups can occur as fundamental groups of smooth complex quasi-projective varieties? The question concerns the abstract group and asks for a full characterization, beyond the supplied partial-family results.
\subsection{Sources}
\begin{itemize}
\item Local UnsolvedMath registry, permanent ID \(30000823\), exact source code \texttt{catalogue entry}, proposed year 2007; no individual proposer is supplied.
\item \emph{Mini-Workshop: Topology of closed one-forms and Cohomology Jumping Loci (2008)}, Oberwolfach Reports, DOI 10.4171/owr/2007/40, pp. 2321--2360. \url{https://doi.org/10.4171/owr/2007/40}.
\item Supplied reference labelled ``graph-manifold characteristic varieties'': \url{https://arxiv.org/abs/1904.00729}.
\item The supplied \texttt{corrected\_verified} note reports checking mathematical content against the official report. The partially solved label and literature assessment checked 2026-08-21 are dated supplied triage requiring independent verification of status and cited theorem scope.
\end{itemize}
\end{document}
